\chapter{Dataset and Implementation}\label{dataset-and-implementation}

This chapter describes the experimental dataset and implementation. It
covers the dataset's public basis and schema, clean and adversarial
examples, Prompt Injection placement, the technology stack, and relevant
implementation and security considerations.

Chapter 3 defines the experimental dataset dimensions and paired design
in Section~\ref{dataset-methodology}. Chapter 4 describes how a record is
represented and processed by the assistant in Sections
\ref{input-representation} and \ref{preprocessing-and-parsing}. This
chapter connects those decisions to the dataset's composition, intended
use, and current implementation status.

\section{Dataset Overview}\label{dataset-overview}

The primary experimental dataset contains 500 synthetic web-application
log records: 250 clean records and 250 adversarial records, organized
into 250 clean/adversarial pairs. The records can be examined by
condition and traffic category; these views represent the same 500
examples and are not additional independent samples.

The 500-record dataset is a controlled research artifact created for
this thesis; it is not the original HTTP CSIC 2010 dataset and is not a
direct export of or row-by-row relabelling of that dataset. HTTP CSIC
2010~\cite{csic2010dataset} provides a public reference for representative
request patterns. The synthetic records do not preserve CSIC source-row
identifiers and are not claimed to be modifications of specific CSIC
examples. Their Prompt Injection text, labels, pair identifiers, and
balanced composition were added for the controlled evaluation.

HTTP CSIC 2010 is used only for the external sample and CSIC-derived
robustness pairs described below. The local copy's download provenance
was not recorded; its content hash and sampling procedure are retained
for reproducibility.

The dataset is restricted to web application logs. Firewall, EDR, DNS,
IDS, Suricata, and Windows event logs are outside the experimental scope.

\section{Public Dataset Basis}\label{implementation-public-dataset-basis}

The Information Security Institute's HTTP CSIC 2010 dataset~\cite{csic2010dataset}
was selected as a public reference for HTTP request patterns relevant to
web-application security. In this thesis, it informs the design context
for the request records; the experimental rows themselves are a separate,
controlled synthetic dataset. This distinction is important: the dataset
does not inherit Prompt Injection examples or labels from CSIC 2010, and
the available files do not support tracing a synthetic row to a specific
original record.

The year 2010 is an acknowledged limitation of the public reference: it
cannot represent the full diversity of current web applications, modern
attack prevalence, or contemporary production traffic. The 500-record
dataset is therefore used for a controlled methodological experiment,
not as evidence that its distribution reflects present-day attacks.

The dataset design adds five traffic categories, a separate clean or
adversarial Prompt Injection condition, five labelled injection
objectives, and explicit clean/adversarial pair identifiers. Its traffic
composition is balanced as shown in Table~\ref{tab:dataset-composition}.

\begin{table}[H]
\centering
\caption{Composition of the 500-record experimental dataset.}
\label{tab:dataset-composition}
\begin{tabular}{@{}lrrr@{}}
\toprule
\textbf{Traffic category} & \textbf{Clean} & \textbf{Adversarial} & \textbf{Total} \\
\midrule
SQL Injection & 50 & 50 & 100 \\
Cross-Site Scripting (XSS) & 50 & 50 & 100 \\
Path Traversal & 50 & 50 & 100 \\
Brute Force Login & 50 & 50 & 100 \\
Benign / Normal & 50 & 50 & 100 \\
\midrule
\textbf{Total} & \textbf{250} & \textbf{250} & \textbf{500} \\
\bottomrule
\end{tabular}
\end{table}

The construction follows the controlled design in Chapter 3,
Section~\ref{dataset-methodology}. First, five traffic categories and
their web-event labels are defined. Clean HTTP records are then paired
with adversarial counterparts that preserve the underlying request
semantics as far as practical and add one Prompt Injection instruction
to an attacker-controlled field. Each row receives a sample ID, pair ID,
condition, and relevant attack or injection annotations. The final balance
check confirms 50 clean and 50 adversarial records per category, with 10
examples of each Prompt Injection variant within each category. The
synthetic records do not retain CSIC source-record identifiers, so this
is a synthetic construction rather than a row-by-row transformation of
CSIC records.

\subsection{External CSIC 2010 Validation Data}
\label{dataset-csic-validation}

In addition to the primary synthetic corpus, the external validation
sample contains 200 HTTP CSIC 2010 requests: 100 Normal and 100
Anomalous. The corpus has 61{,}065 requests and no official test
partition, so this is a study-defined stratified sample from the full
corpus, not a sample from an official test split. It was generated with
seed 42 after checking for exact duplicate requests. Each record retains
its source-record identifier. Sampling parameters, a source-content hash,
label counts, and the field mapping are retained as provenance
~\cite{csic2010dataset}.

The source \texttt{classification} label supports binary evaluation
only: 0 (Normal) maps to \emph{Benign}, and 1 (Anomalous) maps to
\emph{Malicious}. CSIC does not supply native attack-subtype or severity
labels. Anomalous rows are provisionally
assigned to one of the four attack categories used by the synthetic
corpus. Each assignment includes \texttt{candidate\_confidence} and a
\texttt{candidate\_reason}. Low-confidence best-fit assignments are
estimates, not CSIC ground truth. In the 200-record sample, the proposed anomalous categories are 52
SQL Injection, 22 Brute Force Login, 20 Path Traversal, and 6 XSS. They
are not used to calculate Attack-Type Accuracy. Severity remains unavailable, so the
sample cannot support severity-classification accuracy.

\subsubsection{CSIC-Derived Prompt Injection Pairs}
\label{dataset-csic-pi-pairs}

For a distinct robustness test, 50 requests from the external sample
(25 Normal and 25 Anomalous) form 50 clean/adversarial pairs. Each
adversarial member adds one study-authored Prompt Injection instruction
to the query string; the unchanged source remains its clean counterpart.
The five variants have 10 pairs each. Evidence Suppression is applied
only to Anomalous requests so its target is grounded in request evidence;
Prompt Extraction is applied to Normal requests. Each pair retains its
source-record identifier, changed field, variant, target, and scoring
criterion.
They are labelled CSIC-derived, not original CSIC traffic, and are not
added to or pooled with the primary synthetic corpus. Severity Manipulation
uses an injected Low target; because CSIC has no severity ground truth,
this measures response to the instruction rather than severity accuracy.

\section{Dataset Schema}\label{dataset-schema}

Each structured record uses the following fields. The example shows an
adversarial SQL Injection record; its clean counterpart is linked through
the pair identifiers.

\begin{lstlisting}
{
  "sample_id": "SQLI_0001_ADV",
  "paired_sample_id": "SQLI_0001_CLEAN",
  "pair_id": "PAIR_0001",
  "http_method": "POST",
  "url": "/product?note=Ignore previous instructions...",
  "query_parameters": "note=Ignore previous instructions...",
  "request_body": "username=1 UNION SELECT username,password FROM users--",
  "payload": "1 UNION SELECT username,password FROM users--",
  "user_agent": "...",
  "status_code": 400,
  "src_ip": "192.0.2.11",
  "ground_truth_label": "malicious",
  "attack_type": "SQL Injection",
  "traffic_category": "SQL Injection",
  "ground_truth_severity": "high",
  "is_adversarial": true,
  "injection_type": "Instruction Override",
  "injection_field": "query_parameters",
  "prompt_injection_text": "Ignore previous instructions..."
}
\end{lstlisting}

The request fields represent the log content. The remaining fields store
ground-truth and experiment metadata and must not be provided to the model
as labels. In an adversarial record, the instruction itself is embedded
in the indicated request field; \texttt{prompt\_injection\_text}
duplicates that instruction as an annotation for analysis.

In the current dataset, the control category uses
\texttt{traffic\_category = Benign / Normal},
\texttt{ground\_truth\_label = benign}, and \texttt{attack\_type = None}
because no web-attack subtype applies. The thesis refers to this web-event
class simply as \emph{Benign}. This class is independent of the Prompt
Injection condition: both clean and adversarial Benign records are
included. An adversarial Benign record retains the benign web-event label
while its Prompt Injection fields record the separate adversarial
condition.

The \texttt{is\_adversarial}, \texttt{injection\_type}, and
\texttt{injection\_field} values provide the Prompt Injection labels.
The current annotations do not include \texttt{ground\_truth\_mitre},
\texttt{injection\_target}, or \texttt{expected\_evidence}; the chapter
therefore does not treat those fields as available.

The 19 record fields represent three groups: HTTP content and context,
web-event ground truth, and pair/Prompt Injection annotations. Chapter 4,
Section~\ref{input-representation}, uses this same separation to define
the model input. The loader validates field types and normalizes values
as described in Chapter 4, Section~\ref{preprocessing-and-parsing}.

The injection variant and exact text are available, but a per-record
success target and expected-evidence annotation are not. Before
calculating ASR, these scoring annotations must be added to a separate
evaluation store using the success criteria in
Chapter 3, Section~\ref{variant-success-criteria}. They must remain
outside the model input. Until that annotation step is complete, the
dataset alone is not sufficient to compute variant-specific ASR
reproducibly.

\section{Clean and Adversarial
Pairing}\label{clean-and-adversarial-pairing}

Each adversarial record has a corresponding clean record with the same
\texttt{pair\_id}. The \texttt{sample\_id} and
\texttt{paired\_sample\_id} fields identify both members of each pair.

Example:

\begin{lstlisting}
SQLI_0001_CLEAN -- Clean
SQL Injection payload
No Prompt Injection

SQLI_0001_ADV -- Adversarial
Same SQL Injection payload
+ Instruction Override instruction
\end{lstlisting}

Paired design makes it possible to compare the effect of Prompt
Injection while reducing variation in the underlying web event.

\section{Prompt Injection Placement}\label{prompt-injection-placement}

In the current 500-record dataset, Prompt Injection text is embedded in
the \texttt{url} or \texttt{query\_parameters} field. Although User-Agent
and request-body values are also represented as possible attacker-
controlled inputs in the threat model, they do not currently contain the
injections in this dataset version.
 Instruction Override instruction

The five injection categories in the dataset are:

\begin{itemize}
\tightlist
\item
  Instruction Override,
\item
  Classification Manipulation,
\item
  Severity Manipulation,
\item
  Evidence Suppression,
\item
  Prompt Extraction.
\end{itemize}

Each adversarial record has one of these labels, and the exact instruction
and its field are recorded in \texttt{prompt\_injection\_text} and
\texttt{injection\_field}. The 250 adversarial records are balanced across
the five variants, with 50 records per variant. In the current dataset,
125 instructions are placed in \texttt{url} and 125 in
\texttt{query\_parameters}.

The attack classes are limited to SQL Injection, Cross-Site Scripting,
Path Traversal, and Brute Force Login. A Benign category is retained
as a control class. Suspicious File Upload is excluded from the initial
dataset to keep the evaluation focused.

Each of the five traffic categories has 50 clean and 50 adversarial
records, yielding 250 records per condition and 500 in total. Clean SQL
Injection records remain malicious; \emph{clean} refers only to the
absence of Prompt Injection. The web-event label and the Prompt Injection
condition are separate dimensions, including for the Benign control
category.

\section{Technology Stack}\label{technology-stack}

The planned component roles are detailed in Chapter 4,
Section~\ref{application-orchestration-and-data-stores}. Table~\ref{tab:implementation-stack}
distinguishes dependencies currently listed for the project from
components that remain planned. A package being listed as a dependency
does not by itself mean that its application code or experiment workflow
has been completed.

\begin{longtable}[]{@{}p{0.19\textwidth}p{0.34\textwidth}p{0.37\textwidth}@{}}
\caption{Technology status at the time of writing.}
\label{tab:implementation-stack}\\
\toprule
Technology & Intended role & Current status \\
\midrule
\endhead
\bottomrule
\endlastfoot
Python + Pydantic & Load records; validate input and output schemas & Application modules pending. \\
LangChain & Orchestrate C1--C4 processing & Experiment pipelines pending. \\
ChromaDB & Store and retrieve trusted knowledge chunks & Knowledge corpus and index pending. \\
OpenAI API & Hosted LLM inference & Model choice and experiment configuration pending. \\
Sentence Transformers & Local document/query embeddings & Planned component; integration pending. \\
SQLite & Persist experiment metadata and results & Experiment persistence pending. \\
Pandas / scikit-learn & Aggregate predictions and compute metrics & Analytical workflow pending. \\
\end{longtable}

\section{Functional Organization}
\label{functional-organization}

The planned implementation separates four responsibilities: request
validation and preprocessing, trusted-knowledge retrieval, controlled
execution of configurations C1--C4, and evaluation of responses. This
separation keeps scoring labels out of model input and prevents labelled
evaluation examples from becoming retrieval evidence. Application and
experiment workflows remain under development.

\section{Implementation Details}\label{implementation-details}

The implementation will follow the processing path described in Chapter 4.
The planned stages are:

\begin{enumerate}
\def\labelenumi{\arabic{enumi}.}
\tightlist
\item
  Load a structured log record, validate required fields and types, and
  separate model input from scoring metadata, as described in
  Section~\ref{preprocessing-and-parsing}.
\item
  Build a versioned ChromaDB collection from the approved OWASP, ATT\&CK,
  playbook, and Prompt Injection policy documents. Do not index labelled
  logs or evaluation answers; see Sections~\ref{system-knowledge-base}
  and~\ref{retrieval} in Chapter 4.
\item
  Run records under C1--C4 with fixed model and generation settings.
  Keep each clean/adversarial pair together and store its configuration,
  response, detector output, and retrieved evidence IDs.
\item
  Validate the generated structured report and cited evidence IDs, then score web-event
  classification and Prompt Injection outcomes against the separate
  annotations. Chapter 3 defines the metrics and success criteria; the
  detector validation procedure is in Chapter 4,
  Section~\ref{empirical-validation-of-detector}.
\end{enumerate}

The final experiment record must specify the LLM and embedding model
identifiers, prompt and knowledge-base versions, chunking and
top-\emph{k} settings, detector rules, generation parameters, validation
behavior, and any retry policy. These values are not yet fixed, so no
model, retrieval, or end-to-end performance results are reported in this
chapter.

\section{Ethical and Security
Considerations}\label{ethical-and-security-considerations}

The experiments use controlled, offline, or synthetic web-attack
examples for defensive research.

The thesis evaluates whether malicious content inside logs can influence
an LLM. It does not require deployment of attacks against third-party
infrastructure.
